% ============================================================
% ĐỀ ÔN TẬP VẬT LÍ 12 HK1 - THPT TÂN BÌNH
% Chuyển từ PDF sang LaTeX
% Gồm 3 dạng: Trắc nghiệm nhiều lựa chọn - Đúng/Sai - Trả lời ngắn
% Hình vẽ được dựng lại bằng TikZ/PGFPlots
% ============================================================

% ============================================================
% PHẦN I. TRẮC NGHIỆM NHIỀU LỰA CHỌN
% ============================================================

\section*{PHẦN I. TRẮC NGHIỆM NHIỀU LỰA CHỌN}


\nguon{THPT Tân Bình}
\dang{Ôn tập 2}
\begin{ex}
Câu nào sau đây nói về chuyển động của phân tử là \textbf{không đúng}?
\choice
{\True Chuyển động của phân tử là do lực tương tác phân tử gây ra.}
{Các phân tử chuyển động không ngừng.}
{Các phân tử chuyển động càng nhanh thì nhiệt độ càng cao.}
{Các phân tử khí không dao động quanh vị trí cân bằng.}
\loigiai{Chuyển động nhiệt của các phân tử không phải do lực tương tác phân tử gây ra.}
\end{ex}


\nguon{THPT Tân Bình}
\dang{Ôn tập 2}
\begin{ex}
Câu nào sau đây nói về lực tương tác phân tử là \textbf{không đúng}?
\choice
{Lực phân tử chỉ đáng kể khi các phân tử ở rất gần nhau.}
{Lực hút phân tử có thể lớn hơn lực đẩy phân tử.}
{\True Lực hút phân tử không thể lớn hơn lực đẩy phân tử.}
{Lực hút phân tử có thể bằng lực đẩy phân tử.}
\loigiai{Tùy khoảng cách giữa các phân tử, lực hút có thể lớn hơn, bằng hoặc nhỏ hơn lực đẩy.}
\end{ex}


\nguon{THPT Tân Bình}
\dang{Ôn tập 2}
\begin{ex}
Trong các tính chất sau, tính chất nào là của các phân tử chất rắn?
\choice
{Không có hình dạng cố định.}
{Chiếm toàn bộ thể tích của bình chứa.}
{\True Có lực tương tác phân tử lớn.}
{Chuyển động hỗn loạn không ngừng.}
\loigiai{Ở thể rắn, các phân tử liên kết chặt chẽ với nhau nên lực tương tác phân tử lớn.}
\end{ex}


\nguon{THPT Tân Bình}
\dang{Ôn tập 2}
\begin{ex}
Phát biểu nào sau đây nói về chuyển động của phân tử là \textbf{không đúng}?
\choice
{\True Chuyển động của phân tử là do lực tương tác phân tử gây ra.}
{Các phân tử chuyển động không ngừng.}
{Các phân tử chuyển động càng nhanh thì nhiệt độ của vật càng cao.}
{Chuyển động của phân tử là chuyển động nhiệt.}
\loigiai{Chuyển động của phân tử là chuyển động nhiệt, không phải do lực tương tác phân tử gây ra.}
\end{ex}


\nguon{THPT Tân Bình}
\dang{Ôn tập 2}
\begin{ex}
Nguyên tử, phân tử không có tính chất nào sau đây?
\choice
{Chuyển động không ngừng.}
{Giữa chúng có khoảng cách.}
{\True Nở ra khi nhiệt độ tăng, co lại khi nhiệt độ giảm.}
{Chuyển động càng nhanh khi nhiệt độ càng cao.}
\loigiai{Khi nhiệt độ tăng, khoảng cách giữa các hạt có thể tăng làm vật nở ra; bản thân nguyên tử, phân tử không ``nở ra''.}
\end{ex}


\nguon{THPT Tân Bình}
\dang{Ôn tập 2}
\begin{ex}
Yếu tố quyết định vật chất ở thể rắn, thể lỏng hay thể khí là
\choice
{\True lực tương tác giữa các phân tử.}
{màu sắc phân tử.}
{loại phân tử.}
{nhiệt độ khối chất.}
\loigiai{Trạng thái của chất phụ thuộc chủ yếu vào lực tương tác giữa các phân tử và chuyển động nhiệt của chúng.}
\end{ex}


\nguon{THPT Tân Bình}
\dang{Ôn tập 2}
\begin{ex}
Tại sao khi cầm vào vỏ bình ga mini đang sử dụng ta thường thấy có một lớp nước rất mỏng trên đó?
\choice
{Do hơi nước từ tay ta bốc ra.}
{Nước từ trong bình ga thấm ra.}
{\True Do vỏ bình ga lạnh hơn nhiệt độ môi trường nên hơi nước trong không khí ngưng tụ trên đó.}
{Do vỏ bình ga hút nước.}
\loigiai{Khí ga hóa hơi làm vỏ bình lạnh, khiến hơi nước trong không khí ngưng tụ trên bề mặt bình.}
\end{ex}


\nguon{THPT Tân Bình}
\dang{Ôn tập 2}
\begin{ex}
Các tính chất nào sau đây \textbf{không phải} là tính chất của các phân tử chất rắn?
\choice
{Dao động quanh vị trí cân bằng.}
{Lực tương tác phân tử mạnh.}
{Có hình dạng và thể tích xác định.}
{\True Lực tương tác luôn là lực hút.}
\loigiai{Giữa các phân tử chất rắn tồn tại cả lực hút và lực đẩy, tùy khoảng cách.}
\end{ex}


\nguon{THPT Tân Bình}
\dang{Ôn tập 2}
\begin{ex}
Chất rắn vô định hình có đặc tính sau:
\choice
{\True Đẳng hướng và nóng chảy ở nhiệt độ không xác định.}
{Dị hướng và nóng chảy ở nhiệt độ xác định.}
{Dị hướng và nóng chảy ở nhiệt độ không xác định.}
{Đẳng hướng và nóng chảy ở nhiệt độ xác định.}
\loigiai{Chất rắn vô định hình có tính đẳng hướng và không có nhiệt độ nóng chảy xác định.}
\end{ex}


\nguon{THPT Tân Bình}
\dang{Ôn tập 2}
\begin{ex}
Trong các hiện tượng sau đây, hiện tượng nào \textbf{không liên quan} đến sự đông đặc?
\choice
{Tuyết rơi.}
{Đúc tượng đồng.}
{Làm đá trong tủ lạnh.}
{\True Rèn thép trong lò rèn.}
\loigiai{Rèn thép trong lò rèn chủ yếu là quá trình làm nóng và biến dạng kim loại, không phải quá trình đông đặc.}
\end{ex}


\nguon{THPT Tân Bình}
\dang{Ôn tập 2}
\begin{ex}
Chất rắn nào dưới đây là chất rắn vô định hình?
\choice
{\True Thủy tinh.}
{Băng phiến.}
{Hợp kim.}
{Kim loại.}
\loigiai{Thủy tinh là chất rắn vô định hình điển hình.}
\end{ex}


\nguon{THPT Tân Bình}
\dang{Ôn tập 2}
\begin{ex}
Đặc điểm và tính chất nào dưới đây liên quan đến chất rắn vô định hình?
\choice
{Có dạng hình học xác định.}
{Có cấu trúc tinh thể.}
{Có tính dị hướng.}
{\True Không có nhiệt độ nóng chảy xác định.}
\loigiai{Chất rắn vô định hình không có cấu trúc tinh thể và không có nhiệt độ nóng chảy xác định.}
\end{ex}


\nguon{THPT Tân Bình}
\dang{Ôn tập 2}
\begin{ex}
Cho đồ thị biểu diễn sự thay đổi nhiệt độ theo thời gian của nước đá như hình vẽ. Thời gian nước đá tan từ phút nào?

\begin{center}
\begin{tikzpicture}[>=latex,scale=.72]
\draw[->] (0,-1.5) -- (0,3.0) node[above] {Nhiệt độ};
\draw[->] (0,0) -- (11,0) node[right] {Thời gian (phút)};
\foreach \x in {0,1,...,10} \draw[gray!35] (\x,0) -- (\x,2.5);
\foreach \y in {-1,0,1,2} \draw[gray!35] (0,\y) -- (10,\y);
\draw[thick] (0,-1) -- (5,0) -- (7,0) -- (10,2.5);
\draw[dashed] (5,0) -- (5,-1.25) node[below] {6};
\draw[dashed] (7,0) -- (7,-1.25) node[below] {10};
\node[left] at (0,0) {$0$};
\node[left] at (0,-1) {$-5$};
\node[left] at (0,2) {$10$};
\end{tikzpicture}
\end{center}

\choice
{\True Từ phút thứ 6 đến phút thứ 10.}
{Từ phút thứ 10 trở đi.}
{Từ 0 đến phút thứ 6.}
{Từ phút thứ 10 đến phút thứ 15.}
\loigiai{Nước đá tan trong giai đoạn nhiệt độ không đổi. Trên đồ thị, đó là đoạn từ phút thứ 6 đến phút thứ 10.}
\end{ex}


\nguon{THPT Tân Bình}
\dang{Ôn tập 2}
\begin{ex}
Trong suốt thời gian sôi, nhiệt độ của chất lỏng
\choice
{tăng dần lên.}
{giảm dần đi.}
{khi tăng khi giảm.}
{\True không thay đổi.}
\loigiai{Ở điều kiện áp suất xác định, trong suốt quá trình sôi nhiệt độ của chất lỏng không đổi.}
\end{ex}


\nguon{THPT Tân Bình}
\dang{Ôn tập 2}
\begin{ex}
Điều nào sau đây là \textbf{sai} khi nói về sự đông đặc?
\choice
{Sự đông đặc là quá trình chuyển từ thể lỏng sang thể rắn.}
{\True Với một chất rắn, nhiệt độ đông đặc luôn nhỏ hơn nhiệt độ nóng chảy.}
{Trong suốt quá trình đông đặc, nhiệt độ của vật không thay đổi.}
{Nhiệt độ đông đặc của các chất thay đổi theo áp suất bên ngoài.}
\loigiai{Với chất rắn kết tinh, ở cùng một áp suất, nhiệt độ nóng chảy và đông đặc bằng nhau.}
\end{ex}


\nguon{THPT Tân Bình}
\dang{Ôn tập 2}
\begin{ex}
Tốc độ bay hơi của chất lỏng không phụ thuộc vào yếu tố nào sau đây?
\choice
{\True Thể tích của chất lỏng.}
{Gió.}
{Nhiệt độ.}
{Diện tích mặt thoáng của chất lỏng.}
\loigiai{Tốc độ bay hơi phụ thuộc vào nhiệt độ, diện tích mặt thoáng và gió; không phụ thuộc trực tiếp vào thể tích chất lỏng.}
\end{ex}


\nguon{THPT Tân Bình}
\dang{Ôn tập 2}
\begin{ex}
Chọn câu trả lời đúng. Trong sự nóng chảy và đông đặc của các chất rắn:
\choice
{Mỗi chất rắn nóng chảy ở một nhiệt độ xác định, không phụ thuộc vào áp suất bên ngoài.}
{Nhiệt độ đông đặc của chất rắn kết tinh không phụ thuộc áp suất bên ngoài.}
{\True Mỗi chất rắn kết tinh nóng chảy và đông đặc ở cùng một nhiệt độ xác định trong điều kiện áp suất xác định.}
{Mỗi chất rắn nóng chảy ở nhiệt độ nào thì cũng sẽ đông đặc ở nhiệt độ đó.}
\loigiai{Đối với chất rắn kết tinh, ở một áp suất xác định, nhiệt độ nóng chảy và đông đặc là như nhau.}
\end{ex}


\nguon{THPT Tân Bình}
\dang{Ôn tập 2}
\begin{ex}
Gọi $n_R$, $n_L$, $n_K$ lần lượt là mật độ phân tử của một chất ở thể rắn, thể lỏng và thể khí. Thứ tự đúng là
\choice
{$n_R<n_L<n_K$.}
{\True $n_R>n_L>n_K$.}
{$n_R>n_K>n_L$.}
{$n_R<n_K<n_L$.}
\loigiai{Các phân tử ở thể rắn tập trung dày đặc nhất, sau đó đến thể lỏng và thưa nhất là thể khí.}
\end{ex}


% ============================================================
% PHẦN II. TRẮC NGHIỆM ĐÚNG / SAI
% ============================================================

\section*{PHẦN II. TRẮC NGHIỆM ĐÚNG / SAI}


\nguon{THPT Tân Bình}
\dang{Ôn tập 2}
\begin{ex}
Chọn đúng sai khi nói về cấu tạo chất.
\choiceTF
{\True Các chất được cấu tạo từ các hạt riêng gọi là nguyên tử, phân tử.}
{Các nguyên tử, phân tử đứng sát nhau và giữa chúng không có khoảng cách.}
{\True Lực tương tác giữa các phân tử ở thể rắn lớn hơn lực tương tác giữa các phân tử ở thể lỏng và thể khí.}
{\True Các nguyên tử, phân tử chất lỏng dao động xung quanh các vị trí cân bằng không cố định.}
\loigiai{
a) Đúng. \quad
b) Sai, giữa các hạt luôn có khoảng cách. \quad
c) Đúng. \quad
d) Đúng.
}
\end{ex}


\nguon{THPT Tân Bình}
\dang{Ôn tập 2}
\begin{ex}
Cho đồ thị biểu diễn sự thay đổi nhiệt độ theo thời gian của nước đá như hình vẽ.

\begin{center}
\begin{tikzpicture}[>=latex,scale=.75]
\draw[->] (0,-1.8) -- (0,2.4) node[above] {Nhiệt độ ($^\circ$C)};
\draw[->] (0,0) -- (9,0) node[right] {Thời gian (phút)};
\draw[thick] (0,1.6) node[left] {8} -- (3,0) node[below right] {B}
-- (6,0) node[above] {C} -- (9,-1.1) node[below] {D};
\draw[dashed] (3,0) -- (3,-1.4) node[below] {6};
\draw[dashed] (6,0) -- (6,-1.4) node[below] {12};
\draw[dashed] (9,0) -- (9,-1.4) node[below] {18};
\node[above right] at (0,1.6) {A};
\draw[dashed] (0,0) -- (9,0);
\node[left] at (0,0) {$0$};
\end{tikzpicture}
\end{center}

\choiceTF
{Thời gian nước đá đông đặc từ phút thứ 6 đến phút thứ 18.}
{\True Thời gian nước ở thể lỏng từ phút thứ 0 đến phút thứ 6.}
{Thời gian nước ở thể rắn từ phút 0 đến phút thứ 6.}
{\True Thời gian nước đá đông đặc từ phút thứ 6 đến phút thứ 12.}
\loigiai{
a) Sai: đông đặc diễn ra từ phút 6 đến phút 12. \quad
b) Đúng: từ 0 đến 6 phút nước ở thể lỏng và nguội dần. \quad
c) Sai. \quad
d) Đúng.
}
\end{ex}


\nguon{THPT Tân Bình}
\dang{Ôn tập 2}
\begin{ex}
So sánh chất rắn vô định hình và chất rắn kết tinh:
\choiceTF
{\True Khác nhau ở chỗ chất rắn kết tinh có cấu tạo từ những kết cấu rắn có dạng hình học xác định, còn chất rắn vô định hình thì không.}
{Giống nhau ở điểm là cả hai loại chất rắn đều có nhiệt độ nóng chảy xác định.}
{\True Chất rắn kết tinh đa tinh thể có tính đẳng hướng như chất rắn vô định hình.}
{Giống nhau ở điểm cả hai đều có nhiệt độ nóng chảy xác định.}
\loigiai{
a) Đúng. \quad
b) Sai. \quad
c) Đúng. \quad
d) Sai.
}
\end{ex}


\nguon{THPT Tân Bình}
\dang{Ôn tập 2}
\begin{ex}
Hình vẽ đường biểu diễn sự thay đổi nhiệt độ theo thời gian khi đun nóng một chất rắn.

\begin{center}
\begin{tikzpicture}[>=latex,scale=.48]
\draw[->] (0,0) -- (23,0) node[right] {Thời gian (phút)};
\draw[->] (0,0) -- (0,11) node[above] {Nhiệt độ ($^\circ$C)};
\foreach \x in {0,1,...,22} \draw[gray!25] (\x,0)--(\x,10);
\foreach \y in {5,6,...,10} \draw[gray!25] (0,\y)--(22,\y);
\draw[thick] (0,5) -- (5,8) -- (7,8) -- (10,8.8) -- (13,8) -- (18,8) -- (22,7.5);
\draw[dashed] (5,0)--(5,8);
\draw[dashed] (7,0)--(7,8);
\draw[dashed] (13,0)--(13,8);
\draw[dashed] (18,0)--(18,8);
\node[below] at (5,0) {5};
\node[below] at (7,0) {7};
\node[below] at (13,0) {13};
\node[below] at (18,0) {18};
\node[left] at (0,8) {80};
\end{tikzpicture}
\end{center}

\choiceTF
{\True Ở nhiệt độ $80^\circ$C chất rắn này bắt đầu nóng chảy.}
{Thời gian nóng chảy của chất rắn là 4 phút.}
{\True Sự đông đặc bắt đầu vào phút thứ 13.}
{Thời gian đông đặc kéo dài 10 phút.}
\loigiai{
a) Đúng. \quad
b) Sai, thời gian nóng chảy là $7-5=2$ phút. \quad
c) Đúng. \quad
d) Sai, thời gian đông đặc là $18-13=5$ phút.
}
\end{ex}


% ============================================================
% PHẦN III. TRẮC NGHIỆM TRẢ LỜI NGẮN
% ============================================================

\section*{PHẦN III. TRẮC NGHIỆM TRẢ LỜI NGẮN}


\nguon{THPT Tân Bình}
\dang{Ôn tập 2}
\begin{ex}
Cho bảng theo dõi nhiệt độ nóng chảy của chất rắn như sau. Chất rắn bắt đầu nóng chảy phút thứ bao nhiêu?

\begin{center}
\begin{tabular}{|c|c|c|c|c|c|c|}
\hline
\textbf{Thời gian (phút)} & 0 & 2 & 4 & 6 & 8 & 10\\
\hline
\textbf{Nhiệt độ ($^\circ$C)} & 20 & 40 & 60 & 80 & 80 & 85\\
\hline
\end{tabular}
\end{center}

\shortans{6}
\loigiai{Nhiệt độ giữ không đổi ở $80^\circ$C từ phút 6 đến phút 8 nên chất rắn bắt đầu nóng chảy ở phút thứ $6$.}
\end{ex}


\nguon{THPT Tân Bình}
\dang{Ôn tập 2}
\begin{ex}
Hình vẽ đường biểu diễn sự thay đổi nhiệt độ theo thời gian khi đun nóng một chất rắn. Ở nhiệt độ bao nhiêu độ C chất rắn bắt đầu nóng chảy?

\begin{center}
\begin{tikzpicture}[>=latex,scale=.68]
\draw[->] (0,-2.5) -- (0,4.0) node[above] {Nhiệt độ ($^\circ$C)};
\draw[->] (0,0) -- (8.8,0) node[right] {Thời gian (phút)};
\foreach \x in {0,1,...,8} \draw[gray!25] (\x,-2)--(\x,3.2);
\foreach \y in {-2,-1,0,1,2,3} \draw[gray!25] (0,\y)--(8,\y);
\draw[thick] (0,-2) -- (1,0) -- (4,0) -- (7,3);
\node[left] at (0,0) {$0$};
\node[left] at (0,-2) {$-4$};
\node[left] at (0,2) {$4$};
\node[below] at (1,0) {1};
\node[below] at (4,0) {4};
\end{tikzpicture}
\end{center}

\shortans{0}
\loigiai{Chất rắn bắt đầu nóng chảy khi nhiệt độ đạt đến nhiệt độ nóng chảy. Trên đồ thị, nhiệt độ nóng chảy là $0^\circ$C.}
\end{ex}


\nguon{THPT Tân Bình}
\dang{Ôn tập 2}
\begin{ex}
Hình vẽ đường biểu diễn sự thay đổi nhiệt độ theo thời gian khi đun nóng một chất rắn. Thời gian nóng chảy trong bao nhiêu phút?

\begin{center}
\begin{tikzpicture}[>=latex,scale=.68]
\draw[->] (0,-2.5) -- (0,4.0) node[above] {Nhiệt độ ($^\circ$C)};
\draw[->] (0,0) -- (8.8,0) node[right] {Thời gian (phút)};
\foreach \x in {0,1,...,8} \draw[gray!25] (\x,-2)--(\x,3.2);
\foreach \y in {-2,-1,0,1,2,3} \draw[gray!25] (0,\y)--(8,\y);
\draw[thick] (0,-2) -- (1,0) -- (4,0) -- (7,3);
\node[left] at (0,0) {$0$};
\node[below] at (1,0) {1};
\node[below] at (4,0) {4};
\end{tikzpicture}
\end{center}

\shortans{3}
\loigiai{Quá trình nóng chảy diễn ra trong đoạn nhiệt độ không đổi từ phút 1 đến phút 4:
\[
\Delta t=4-1=3\text{ phút}.
\]
}
\end{ex}


\nguon{THPT Tân Bình}
\dang{Ôn tập 2}
\begin{ex}
Một chất có thể có tối đa bao nhiêu sự chuyển thể?
\shortans{3}
\loigiai{Có ba dạng chuyển thể cơ bản: nóng chảy (rắn $\to$ lỏng), hóa hơi (lỏng $\to$ khí) và thăng hoa (rắn $\to$ khí).}
\end{ex}


\nguon{THPT Tân Bình}
\dang{Ôn tập 2}
\begin{ex}
Dựa vào đồ thị sự thay đổi nhiệt độ của chất kết tinh được làm nóng chảy. Em hãy cho biết chất rắn nóng chảy hoàn toàn sau bao nhiêu giây?

\begin{center}
\begin{tikzpicture}[>=latex,scale=.75]
\draw[->] (0,0) -- (8.5,0) node[right] {Thời gian (s)};
\draw[->] (0,0) -- (0,4.2) node[above] {Nhiệt độ ($^\circ$C)};
\draw[thick] (0,0) -- (2,2.0) -- (6,2.0) -- (8,3.5);
\draw[dashed] (2,0)--(2,2);
\draw[dashed] (6,0)--(6,2);
\node[below] at (2,0) {20};
\node[below] at (6,0) {60};
\node[below] at (8,0) {80};
\node[left] at (0,2) {$t_c$};
\node[above] at (4,2) {nóng chảy};
\end{tikzpicture}
\end{center}

\shortans{60}
\loigiai{Chất rắn nóng chảy hoàn toàn khi kết thúc đoạn nhiệt độ không đổi, tức tại $t=60$ s.}
\end{ex}


\nguon{THPT Tân Bình}
\dang{Ôn tập 2}
\begin{ex}
Dựa vào đồ thị sự thay đổi nhiệt độ của nước theo thời gian khi được đun sôi. Hãy cho biết thời gian nước sôi là bao lâu?

\begin{center}
\begin{tikzpicture}[>=latex,scale=.75]
\draw[->] (0,0) -- (8.5,0) node[right] {Thời gian (s)};
\draw[->] (0,0) -- (0,4.0) node[above] {Nhiệt độ ($^\circ$C)};
\draw[thick] (0,0.8) node[left] {20} -- (3.6,3.0) node[above] {B}
-- (7.8,3.0) node[above] {C};
\draw[dashed] (3.6,0)--(3.6,3.0);
\draw[dashed] (7.8,0)--(7.8,3.0);
\node[below] at (3.6,0) {60};
\node[below] at (7.8,0) {130};
\node[left] at (0,3.0) {$100^\circ$C};
\end{tikzpicture}
\end{center}

\shortans{70}
\loigiai{Nước sôi từ $t=60$ s đến $t=130$ s nên:
\[
\Delta t=130-60=70\text{ s}.
\]
}
\end{ex}


% ============================================================
% ĐÁP ÁN TỔNG HỢP
% ============================================================

\section*{ĐÁP ÁN}

\textbf{Phần I:}
\[
1A,\ 2C,\ 3C,\ 4A,\ 5C,\ 6A,\ 7C,\ 8D,\ 9A,\ 10D,
\]
\[
11A,\ 12D,\ 13A,\ 14D,\ 15B,\ 16A,\ 17C,\ 18B.
\]

\textbf{Phần II:}
\[
\begin{array}{c|cccc}
\text{Câu} & a & b & c & d\\ \hline
1 & Đ & S & Đ & Đ\\
2 & S & Đ & S & Đ\\
3 & Đ & S & Đ & S\\
4 & Đ & S & Đ & S
\end{array}
\]

\textbf{Phần III:}
\[
1.\ 6;\qquad
2.\ 0;\qquad
3.\ 3;\qquad
4.\ 3;\qquad
5.\ 60;\qquad
6.\ 70.
\]
